\documentclass[aspectratio=169]{beamer}

% Theme and colors to match the blue/gray style
\usetheme{default}
\usecolortheme{default}

% Custom colors
\definecolor{slidebg}{RGB}{55,55,60}
\definecolor{bluebar}{RGB}{60,80,120}
\definecolor{darkterm}{RGB}{25,25,25}
\definecolor{termgreen}{RGB}{100,220,100}
\definecolor{termwhite}{RGB}{220,220,220}
\definecolor{termyellow}{RGB}{220,220,100}
\definecolor{termred}{RGB}{220,80,80}
\definecolor{termcyan}{RGB}{100,220,220}
\definecolor{highlight}{RGB}{200,0,0}
\definecolor{arrowred}{RGB}{220,30,30}

% White background, standard dark text
\setbeamercolor{frametitle}{bg=bluebar,fg=black}
\setbeamercolor{normal text}{fg=black}
\setbeamercolor{structure}{fg=bluebar}

% Remove navigation symbols
\setbeamertemplate{navigation symbols}{}

% Frame title template with blue bar
\setbeamertemplate{frametitle}{
  \vspace{-0.1cm}
  \begin{beamercolorbox}[wd=\paperwidth,ht=1.2cm,dp=0.3cm]{frametitle}
    \hspace{0.5cm}\usebeamerfont{frametitle}\insertframetitle
  \end{beamercolorbox}
}

\setbeamerfont{frametitle}{size=\LARGE}

% Packages
\usepackage[T1]{fontenc}
\usepackage{helvet}                    % Helvetica — matches the original slide font
\renewcommand{\familydefault}{\sfdefault}
\usepackage{listings}
\usepackage{textcomp}
\usepackage{tikz}
% Base listing settings (no backgroundcolor — tcolorbox handles the box)
\lstset{
  basicstyle=\ttfamily\scriptsize\color{termwhite},
  breaklines=true,
  showstringspaces=false,
  upquote=true,
  columns=fullflexible,
  aboveskip=0pt,
  belowskip=0pt,
}

\lstdefinestyle{terminal}{ basicstyle=\ttfamily\scriptsize\color{termwhite} }
\lstdefinestyle{terminalsmall}{ basicstyle=\ttfamily\tiny\color{termwhite} }
\lstdefinestyle{config}{ basicstyle=\ttfamily\small\color{termwhite} }

\usepackage{tcolorbox}
\tcbuselibrary{listings,skins}

% tcolorbox wraps the listing in one continuous colored box (no per-line stripes)
\newtcblisting{term}{
  colback=darkterm, colframe=darkterm,
  arc=4pt, boxrule=0pt,
  left=6pt, right=6pt, top=4pt, bottom=4pt,
  listing only,
  listing options={style=terminal},
}

\newtcblisting{termsmall}{
  colback=darkterm, colframe=darkterm,
  arc=4pt, boxrule=0pt,
  left=6pt, right=6pt, top=4pt, bottom=4pt,
  listing only,
  listing options={style=terminalsmall},
}

\newtcblisting{conf}{
  colback=darkterm, colframe=darkterm,
  arc=4pt, boxrule=0pt,
  left=6pt, right=6pt, top=4pt, bottom=4pt,
  listing only,
  listing options={style=config},
}

\title{\Huge Secure DevOps}
\subtitle{\Large Setting up a secure environment for Ansible}
\author{}
\date{}

\begin{document}

% ==============================================================
% Slide 1 - Title
% ==============================================================
\begin{frame}
\vspace{2cm}
\begin{beamercolorbox}[wd=\paperwidth,ht=3cm,dp=1cm]{frametitle}
  \hspace{1cm}{\Huge\textbf{Secure DevOps}}\\[0.3cm]
  \hspace{1cm}{\Large Setting up a secure environment for Ansible}
\end{beamercolorbox}
\vfill
\hfill\begin{minipage}{0.3\textwidth}
\raggedleft
{\small PRESENTED BY:}\\
{\Large\textbf{Joel Kirch}}
\end{minipage}\hspace{1cm}
\vspace{0.5cm}

\begin{center}
{\tiny\textcopyright\ CC BY-NC-ND --- This license enables reusers to copy and distribute the material in any medium or format in unadapted form only, for noncommercial purposes only, and only so long as attribution is given to the creator.}
\end{center}
\end{frame}

% ==============================================================
% Slide 2 - Today's Topics
% ==============================================================
\begin{frame}{Today's Topics}
\begin{itemize}
  \item Many demos and tutorials skip over how to set up a secure environment for Ansible
  \begin{itemize}
    \item SSH \& GPG Keys, SSH-Agent, GPG-Agent
    \item \texttt{pass} -- The Unix Password Manager
    \item Ansible Vault
  \end{itemize}
  \item \textbf{SSH Setup Demo}
  \item \textbf{Ansible Setup Demo}
  \item \textbf{Ansible Demo}
\end{itemize}
\end{frame}

% ==============================================================
% Slide 3 - Section Divider: SSH Setup Demo
% ==============================================================
\begin{frame}[plain]
\vspace{3cm}
\begin{beamercolorbox}[wd=\paperwidth,ht=2.5cm,dp=0.8cm]{frametitle}
  \centering{\Huge SSH Setup Demo}
\end{beamercolorbox}
\end{frame}

% ==============================================================
% Slide 4 - Generate SSH Keys
% ==============================================================
\begin{frame}[fragile]{Generate SSH Keys}
\begin{itemize}
  \item Generate an ED25519 SSH key with the comment ``DemoSSHKey'' and save it using that same filename -- \textcolor{highlight}{USE A PASSWORD to protect your SSH Key.}
\end{itemize}
\vspace{0.2cm}
\begin{conf}
ssh-keygen -t ed25519 -C "DemoSSHKey" -f DemoSSHKey
\end{conf}
\begin{itemize}
  \item This command breaks down as follows:
  \begin{itemize}
    \item \textbf{\texttt{-t ed25519}}: Specifies the type of key to generate (ED25519).
    \item \textbf{\texttt{-C "DemoSSHKey"}}: Adds the comment ``DemoSSHKey'' to the key
    \item \textbf{\texttt{-f DemoSSHKey}}: Saves the private key to a file named \texttt{DemoSSHKey} and the public key to a file named \texttt{DemoSSHKey.pub}.
  \end{itemize}
  \item After running this command, you'll have two files:
  \begin{itemize}
    \item DemoSSHKey (the private key)
    \item DemoSSHKey.pub (the corresponding public key)
  \end{itemize}
\end{itemize}
\end{frame}

% ==============================================================
% Slide 5 - Upload SSH Public Key
% ==============================================================
\begin{frame}[fragile]{Upload SSH Public Key}
\begin{itemize}
  \item Copy your public SSH key and upload it to your server or cloud provider
  \item The public key is safe to share -- that's the whole point
  \item For cloud providers (e.g., DigitalOcean), paste the key in the SSH Keys section of the console
\end{itemize}
\vspace{0.3cm}
\begin{conf}
cat ~/.ssh/DemoSSHKey.pub
\end{conf}
\end{frame}

% ==============================================================
% Slide 6 - RSA 4096 vs ED25519
% ==============================================================
\begin{frame}[fragile]{RSA 4096 vs ED25519}
\begin{termsmall}
ssh-rsa AAAAB3NzaC1yc2EAAAADAQABAAACAQD4kiiiQH2Avx4rzOsNf
fLvnEZAMuEF7CA0cpg3L6BO1aFu6jEFvr/NonnL1tPN/OwDoTeEMG4k7B
N3dAXFFkgFAGAT+2RNvTnFb+nhyReeEkwZg+Iw9+stfxITrtvKgoNQ76n
884P79FciMbn9QP6BmCXsTUQhMFMvWHP03rwYXFt9f/pUgW6QCmpIw5i3f
KgQf0WGFVsjPWXMCXrNiIDDk/71VUqXg8FKFUvQUG1DGDcspPdel0gsJSX
y14Wov4l4sSlViYqe6WQXKA6aIvsQGk+DfZErKPWMg+20LCAYYQeD1+r4Q
iLbMHbHNK26WBmkjzv0pgkbEdBdrscRaTi8AFCWelyqsoxxz4GEK0CG7NoB
idl0KTf95OqDpmAnZcwRDTLWbYlLBxgL8+Pmjgku9Pnzhf2qipQs94A+te
EgZiPlpIXSitqcBmJ0aamCTMrYE9BkylM7hvGkfx7QZYWktQkgdPKKl8ac
KDM77nc1SWc1vQ2pHhq53xk08cNWrOMZaxqZaDMGncxp83GnOEjkZhr6Xw
9T5/ Demo_SSH_RSA_Key
\end{termsmall}
\vspace{0.2cm}
\begin{term}
ssh-ed25519 AAAAC3NzaC1lZDI1NTE5AAAAIDGwqTc2eyAVxspuLyNERj
9ktEHdkRwdvKpAaO8ZP68 Demo_SSH_ED25519_Key
\end{term}
\end{frame}

% ==============================================================
% Slide 7 - Configure SSH Client
% ==============================================================
\begin{frame}[fragile]{Configure SSH Client}
\begin{conf}
Host ansible
    HostName 192.168.122.100
    User root
    IdentityFile ~/.ssh/DemoSSHKey
\end{conf}
\vspace{0.3cm}
\textbf{\texttt{Host ansible}}: Creates an alias ``ansible'' for easy connection.\\[0.2cm]
\textbf{\texttt{HostName 192.168.122.100}}: Sets the server's IP address.\\[0.2cm]
\textbf{\texttt{User root}}: Uses ``root'' as the default username.\\[0.2cm]
\textbf{\texttt{IdentityFile \textasciitilde/.ssh/DemoSSHKey}}: Specifies the SSH key for secure login.
\end{frame}

% ==============================================================
% Slide 8 - SSH to ansible, part 1
% ==============================================================
\begin{frame}[fragile]{SSH to ansible, part 1}
\begin{columns}[T]
\begin{column}{0.55\textwidth}
\begin{termsmall}
$ ssh ansible
Enter passphrase for key '~/.ssh/DemoSSHKey':
Welcome to Ubuntu 24.04 LTS

 System load:  0.08     Processes:        104
 Usage of /:   19.2%    Users logged in:  0
 Memory usage: 36%      IPv4: 192.168.122.100

root@ansible:~#
\end{termsmall}
\end{column}
\begin{column}{0.42\textwidth}
\begin{itemize}
  \item Notice that I have to enter a password to use my SSH key.
  \item Let's fix that by using SSH agent
  \item Add this line to your config file:\\
    \textbf{\texttt{AddKeysToAgent yes}}
\end{itemize}
\vspace{0.2cm}
\textcolor{arrowred}{$\Rightarrow$}
\begin{termsmall}
Host ansible
    HostName 192.168.122.100
    User root
    AddKeysToAgent yes
    IdentityFile ~/.ssh/DemoSSHKey
\end{termsmall}
\end{column}
\end{columns}
\end{frame}

% ==============================================================
% Slide 9 - CyberSecurity Time
% ==============================================================
\begin{frame}[fragile]{CyberSecurity Time}
\begin{columns}[T]
\begin{column}{0.55\textwidth}
\begin{itemize}
  \item Update the \texttt{sshd\_config} on the ansible server
  \begin{itemize}
    \item Disable root login: \textbf{\texttt{PermitRootLogin no}} ensures that root cannot access the server via SSH.
    \item Enforce public key authentication: \textbf{\texttt{PubkeyAuthentication yes}} and \textbf{\texttt{PasswordAuthentication no}} require that all users authenticate with an SSH key.
    \item Do not forget to restart SSHD: \textbf{\texttt{systemctl restart ssh.service}}
  \end{itemize}
\end{itemize}
\end{column}
\begin{column}{0.42\textwidth}
\begin{termsmall}
# Authentication:

#LoginGraceTime 2m
PermitRootLogin no          <--
#StrictModes yes
#MaxAuthTries 6

PubkeyAuthentication yes    <--

...

PasswordAuthentication no   <--
#PermitEmptyPasswords no
\end{termsmall}
\end{column}
\end{columns}
\end{frame}

% ==============================================================
% Slide 10 - Update ssh config
% ==============================================================
\begin{frame}[fragile]{Update ssh config}
\begin{columns}[T]
\begin{column}{0.45\textwidth}
\begin{itemize}
  \item Update my ssh client config to use \textbf{\texttt{ansible\_user}} instead of \textbf{\texttt{root}}
\end{itemize}
\vspace{0.3cm}
\begin{term}
Host ansible
    HostName 192.168.122.100
    User ansible_user
    AddKeysToAgent yes
    IdentityFile ~/.ssh/DemoSSHKey
\end{term}
\end{column}
\begin{column}{0.52\textwidth}
\begin{termsmall}
$ ssh ansible
root@192.168.122.100: Permission denied
  (publickey).

$ vim ~/.ssh/config

$ ssh ansible
Welcome to Ubuntu 24.04 LTS
...
ansible_user@ansible:~$
\end{termsmall}
\end{column}
\end{columns}
\end{frame}

% ==============================================================
% Slide 11 - SSH Setup Summary
% ==============================================================
\begin{frame}{SSH Setup Summary}
\begin{itemize}
  \item Created SSH keys, uploaded the public key to cloud provider
  \item Configured our client to use SSH Agent with SSH keys
  \item Connected to our ansible server, using SSH keys
  \item Hardened SSHd on the ansible server
  \item Verified that our settings worked
  \item \textbf{QUESTION: Are we doing DevOps?}
\end{itemize}
\end{frame}

% ==============================================================
% Slide 12 - Section Divider: Ansible Setup Demo
% ==============================================================
\begin{frame}[plain]
\vspace{3cm}
\begin{beamercolorbox}[wd=\paperwidth,ht=2.5cm,dp=0.8cm]{frametitle}
  \centering{\Huge Ansible Setup Demo}
\end{beamercolorbox}
\end{frame}

% ==============================================================
% Slide 13 - Install Ansible
% ==============================================================
\begin{frame}[fragile]{Install Ansible}
\begin{termsmall}
ansible_user@ansible:~$ ./install_ansible.sh
Reading package lists... Done
Building dependency tree... Done
...
The following NEW packages will be installed:
  python3 python3-pip python3-venv ...
...
Do you want to continue? [Y/n]
\end{termsmall}
\vspace{0.3cm}
\begin{itemize}
  \item Use \texttt{git} to download the install script
  \item \texttt{chmod +x install\_ansible.sh}
  \item \texttt{./install\_ansible.sh}
\end{itemize}
\end{frame}

% ==============================================================
% Slide 14 - Logout & Login
% ==============================================================
\begin{frame}[fragile]{Logout \& Login}
\begin{termsmall}
ansible_user@ansible:~$ ./install_ansible.sh
...
~/.local/bin has been added to PATH, but you need to open
a new terminal or re-login for this PATH change to take effect.

ansible_user@ansible:~$ ansible --version
ansible [core 2.18.1]
  config file = None
  ...
  python version = 3.12.3
\end{termsmall}
\end{frame}

% ==============================================================
% Slide 15 - Ansible Installed
% ==============================================================
\begin{frame}[fragile]{Ansible Installed}
\begin{term}
ansible_user@ansible:~$ ansible --version
ansible [core 2.18.1]
  config file = None
  configured module search path =
    ['/home/ansible_user/.ansible/plugins/modules',
     '/usr/share/ansible/plugins/modules']
  ansible python module location =
    /home/ansible_user/.local/share/uv/...
  executable location =
    /home/ansible_user/.local/bin/ansible
  python version = 3.12.3
  jinja version = 3.1.4
  libyaml = True
\end{term}
\end{frame}

% ==============================================================
% Slide 16 - Section Divider: Ansible Demo
% ==============================================================
\begin{frame}[plain]
\vspace{3cm}
\begin{beamercolorbox}[wd=\paperwidth,ht=2.5cm,dp=0.8cm]{frametitle}
  \centering{\Huge Ansible Demo}
\end{beamercolorbox}
\end{frame}

% ==============================================================
% Slide 17 - Our First Ansible Playbook
% ==============================================================
\begin{frame}[fragile]{Our First Ansible Playbook}
\begin{columns}[T]
\begin{column}{0.35\textwidth}
\begin{itemize}
  \item Create an inventory file
  \item Write a playbook
  \item Run the playbook
\end{itemize}
\end{column}
\begin{column}{0.62\textwidth}
\begin{termsmall}
[local]
localhost ansible_connection=local
\end{termsmall}
\vspace{0.1cm}
\begin{termsmall}
---
- name: Ping localhost
  hosts: local
  tasks:
    - name: Ping the localhost
      ping:
\end{termsmall}
\end{column}
\end{columns}
\vspace{0.2cm}
\begin{termsmall}
ansible_user@ansible:~/secure_devops_demo$ ansible-playbook ping.yml

PLAY [Ping localhost] *****************************************************
TASK [Gathering Facts] ****************************************************
ok: [localhost]
TASK [Ping the localhost] *************************************************
ok: [localhost]
PLAY RECAP ****************************************************************
localhost           : ok=2    changed=0    unreachable=0    failed=0
\end{termsmall}
\end{frame}

% ==============================================================
% Slide 18 - ansible.cfg
% ==============================================================
\begin{frame}[fragile]{ansible.cfg}
\begin{columns}[T]
\begin{column}{0.45\textwidth}
\begin{term}
[defaults]
inventory = inventory.ini
command_warnings = False
deprecation_warnings = False
interpreter_python =
    /usr/bin/python3
\end{term}
\vspace{0.3cm}
Add the inventory location and turn off those warnings
\end{column}
\begin{column}{0.52\textwidth}
Now you don't have to pass the \texttt{-i inventory.ini} each time you run a playbook
\vspace{0.3cm}
\begin{termsmall}
ansible_user@ansible:~/secure_devops_demo$
  ansible-playbook ping.yml

PLAY [Ping localhost] *****************
TASK [Gathering Facts] ****************
ok: [localhost]
TASK [Ping the localhost] *************
ok: [localhost]
PLAY RECAP ****************************
localhost  : ok=2  changed=0  failed=0
\end{termsmall}
\end{column}
\end{columns}
\end{frame}

% ==============================================================
% Slide 19 - Update ansible with Ansible
% ==============================================================
\begin{frame}[fragile]{Update ansible with Ansible}
\begin{columns}[T]
\begin{column}{0.48\textwidth}
\begin{term}
---
- name: Update and upgrade packages
    on localhost
  hosts: local
  become: yes  tasks:
    - name: Update apt cache
      apt:
        update_cache: yes
    - name: Upgrade all packages
      apt:
        upgrade: dist
\end{term}
\end{column}
\begin{column}{0.48\textwidth}
\begin{itemize}
  \item Notice that we are using \textbf{\texttt{become}} to use sudo
  \item It failed?
  \item Why?
\end{itemize}
\end{column}
\end{columns}
\vspace{0.2cm}
\begin{termsmall}
ansible_user@ansible:~/secure_devops_demo$ ansible-playbook update.yml
PLAY [Update and upgrade packages on localhost] ***************************
TASK [Gathering Facts] ****************************************************
fatal: [localhost]: FAILED! => {"msg": "sudo: a password is required\n"}
\end{termsmall}
\end{frame}

% ==============================================================
% Slide 20 - Ansible with sudo
% ==============================================================
\begin{frame}[fragile]{Ansible with sudo}
\begin{term}
ansible_user@ansible:~/secure_devops_demo$
  ansible-playbook update.yml --ask-become-pass
BECOME password:

PLAY [Update and upgrade packages on localhost] ********
TASK [Gathering Facts] *********************************
ok: [localhost]
TASK [Update apt cache] ********************************
changed: [localhost]
\end{term}
\vspace{0.3cm}
\begin{itemize}
  \item It worked!
  \item Typing passwords each time you run a playbook is no good.
  \item Let's fix that.
\end{itemize}
\end{frame}

% ==============================================================
% Slide 21 - Storing Secrets in Ansible
% ==============================================================
\begin{frame}{Storing Secrets in Ansible}
\begin{itemize}
  \item Ansible can store secrets in a ``vault''
  \begin{itemize}
    \item Uses AES 256 to encrypt the data
    \item Needs to be unlocked before using
  \end{itemize}
  \vspace{0.3cm}
  \item You could store the \texttt{ansible\_user} password for sudo in the vault and run the playbook:\\
    \textbf{\texttt{ansible-playbook update.yml --ask-vault-pass}}
  \vspace{0.3cm}
  \item \textbf{Not much better than just running:}\\
    \textbf{\texttt{ansible-playbook update.yml --ask-become-pass}}
\end{itemize}
\end{frame}

% ==============================================================
% Slide 22 - Ansible Vault
% ==============================================================
\begin{frame}[fragile]{Ansible Vault}
\begin{term}
ansible_user@ansible:~/secure_devops_demo$ ansible-vault view vault.yml
Vault password:
ansible_become_password: ************
\end{term}
\vspace{0.3cm}
\begin{termsmall}
ansible_user@ansible:~/secure_devops_demo$ cat vault.yml
$ANSIBLE_VAULT;1.1;AES256
61366165396139306136663666343065626638343333366264393063643130
33663663363633363634323361313435393836323133303933343631333830
36306533313633393136343436313634333536333533333336343939363636
38663962303633363939306133383631333836333738333330363633373631
\end{termsmall}
\end{frame}

% ==============================================================
% Slide 23 - Lets use pass
% ==============================================================
\begin{frame}[fragile]{Let's use \texttt{pass}}
\begin{termsmall}
ansible_user@ansible:~/secure_devops_demo$ ./install_setup_pass.sh
Installing pass, the Unix password manager...
[sudo] password for ansible_user:

gpg: directory '/home/ansible_user/.gnupg' created
gpg: keybox '/home/ansible_user/.gnupg/pubring.kbx' created
Generating a new GPG key for pass...
...
Password store initialized for 2466552D8747919
Pass initialized with GPG key ID: 2466552D8747919
Enter the password for the Ansible Vault:

Enter contents of ansible/vault_password and press Ctrl+D when finished:

Ansible Vault password has been securely stored in pass under
  'ansible/vault_password'.
ansible_user@ansible:~/secure_devops_demo$ pass ansible/vault_password
\end{termsmall}
\end{frame}

% ==============================================================
% Slide 24 - pass uses GPG
% ==============================================================
\begin{frame}[fragile]{\texttt{pass} uses GPG}
\centering{\large \texttt{pass} uses GPG to encrypt/decrypt passwords}
\vspace{0.2cm}
\begin{columns}[T]
\begin{column}{0.48\textwidth}
\begin{termsmall}
ansible_user@ansible:~$ tree
  ~/.password-store/
/home/ansible_user/.password-store/
  ansible
    vault_password.gpg
\end{termsmall}
\vspace{0.2cm}
\texttt{pass} just stores each entry as its own GPG encrypted file
\vspace{0.3cm}

GPG uses an agent, just like the SSH Agent
\begin{termsmall}
/run/user/1000/gnupg/S.gpg-agent

SSH Agent PID:
25296 ssh-agent
25299 ssh-agent
\end{termsmall}
\end{column}
\begin{column}{0.48\textwidth}
\begin{termsmall}
ansible_user@ansible:~$ file
  ~/.password-store/ansible/
    vault_password.gpg
...PGP RSA encrypted session key
  - keyid: DEF15331 A3CB758D
  RSA (Encrypt or Sign) 2048b

ansible_user@ansible:~$ xxd
  ~/.password-store/ansible/
    vault_password.gpg
00000000: 8501 0c03 def1 5331
00000010: 8ee9 9d32 2e5a 0c68
00000020: bf34 3e73 2e26 32b6
...
\end{termsmall}
\end{column}
\end{columns}
\end{frame}

% ==============================================================
% Slide 25 - Update ansible.cfg for pass
% ==============================================================
\begin{frame}[fragile]{Update \texttt{ansible.cfg} for \texttt{pass}}
\begin{itemize}
  \item Update the ansible.cfg to call a script to use pass to unlock our vault
\end{itemize}
\vspace{0.2cm}
\begin{columns}[T]
\begin{column}{0.55\textwidth}
\begin{term}
[defaults]
inventory = inventory.ini
command_warnings = False
deprecation_warnings = False
interpreter_python = /usr/bin/python3
vault_password_file =
    bin/get_vault_pass.sh
\end{term}
\end{column}
\begin{column}{0.4\textwidth}
\begin{term}
#!/bin/bash
pass ansible/vault_password
\end{term}
\end{column}
\end{columns}
\vspace{0.3cm}
\begin{itemize}
  \item Wait, what were we doing again?
  \item Updating the packages on the ansible server, securely using \textcolor{highlight}{NO PLAINTEXT passwords}.
\end{itemize}
\end{frame}

% ==============================================================
% Slide 26 - Update our update.yml
% ==============================================================
\begin{frame}[fragile]{Update our \texttt{update.yml}}
\begin{columns}[T]
\begin{column}{0.48\textwidth}
\begin{term}
---
- name: Update and upgrade packages
    on localhost
  hosts: local
  become: yes
  vars_files:    - vault.yml  tasks:
    - name: Update apt cache
      apt:
        update_cache: yes
    - name: Upgrade all packages
      apt:
        upgrade: dist
\end{term}
\end{column}
\begin{column}{0.48\textwidth}
\begin{termsmall}
ansible_user@ansible:~/secure_devops_demo$
  ansible-playbook update.yml

PLAY [Update and upgrade packages
  on localhost] ********************

TASK [Gathering Facts] *************
ok: [localhost]

TASK [Update apt cache] ************
changed: [localhost]

TASK [Upgrade all packages] ********
ok: [localhost]

PLAY RECAP *************************
localhost  : ok=3  changed=1
  unreachable=0  failed=0
\end{termsmall}
\end{column}
\end{columns}
\end{frame}

% ==============================================================
% Slide 27 - Okay...
% ==============================================================
\begin{frame}{Okay...}
\begin{itemize}
  \item Couldn't I have just run:\\[0.2cm]
    \texttt{apt update \&\& apt upgrade -y}
  \vspace{0.5cm}
  \item \textbf{Yes, but that's not DevOps}
  \begin{itemize}
    \item It doesn't scale to dozens / hundreds / thousands of machines
    \item Error prone
    \item You don't get ``infrastructure as code''
  \end{itemize}
\end{itemize}
\end{frame}

% ==============================================================
% Slide 28 - Let's spin up more machines
% ==============================================================
\begin{frame}[fragile]{Let's spin up more machines}
\begin{itemize}
  \item 3 local VMs on QEMU/KVM -- no internet needed
  \item Add them to \textbf{\texttt{inventory.ini}} so Ansible can reach them
\end{itemize}
\vspace{0.3cm}
\begin{term}
ansible_user@ansible:~/secure_devops_demo$ bin/update_local_inventory.sh
Local VM IPs (static):
  target-1: 192.168.122.101
  target-2: 192.168.122.102
  target-3: 192.168.122.103
Updated inventory.ini
Cleared stale known_hosts entries for target IPs
Done. Run bin/update_ssh_config.sh after creating ansible_user.
\end{term}
\end{frame}

% ==============================================================
% Slide 29 - Ansible to update .bashrc
% ==============================================================
\begin{frame}[fragile]{Ansible to update \texttt{.bashrc}}
\begin{columns}[T]
\begin{column}{0.58\textwidth}
\begin{termsmall}
---
- hosts: localhost
  tasks:
    - name: Ensure ssh-agent is started
        and environment variables are saved
      blockinfile:
        path: ~/.bashrc
        marker: "# {mark} Ansible SSH
          agent configuration"
        block: |
          # Ansible modified: SSH agent setup
          if [ -z "$SSH_AUTH_SOCK" ]; then
            eval "$(ssh-agent -s)" >
              ~/.ssh-agent-variables
            ssh-add ~/.ssh/DemoSSHKey
          fi
          if [ -f ~/.ssh-agent-variables ]
          then
            source ~/.ssh-agent-variables
          fi
        create: yes
\end{termsmall}
\end{column}
\begin{column}{0.38\textwidth}
Now the \texttt{.bashrc} file is updated to start the SSH-Agent when the user logs in
\vspace{0.3cm}
\begin{termsmall}
# BEGIN Ansible SSH agent
#   configuration
# Ansible modified: SSH agent
#   setup
if [ -z "$SSH_AUTH_SOCK" ]; then
    eval "$(ssh-agent -s)" >
      ~/.ssh-agent-variables
    ssh-add ~/.ssh/DemoSSHKey
fi
if [ -f ~/.ssh-agent-variables ]
then
    source ~/.ssh-agent-variables
fi
# END Ansible SSH agent
#   configuration
\end{termsmall}
\end{column}
\end{columns}
\end{frame}

% ==============================================================
% Slide 30 - Idempotent
% ==============================================================
\begin{frame}[fragile]{Idempotent}
\begin{columns}[T]
\begin{column}{0.55\textwidth}
\begin{termsmall}
ansible_user@ansible:~/secure_devops_demo$
  ansible-playbook add_ssh_key.yml

PLAY [localhost] ***********************
TASK [Gathering Facts] *****************
ok: [localhost]
TASK [Ensure ssh-agent is started...] **
changed: [localhost]

PLAY RECAP *****************************
localhost  : ok=2  changed=1

ansible_user@ansible:~/secure_devops_demo$
  ansible-playbook add_ssh_key.yml

PLAY [localhost] ***********************
TASK [Gathering Facts] *****************
ok: [localhost]
TASK [Ensure ssh-agent is started...] **
ok: [localhost]

PLAY RECAP *****************************
localhost  : ok=2  changed=0
\end{termsmall}
\end{column}
\begin{column}{0.42\textwidth}
\begin{itemize}
  \item Idempotency ensures repeatability: Running the same operation multiple times produces the same outcome without unintended side effects.
  \vspace{0.3cm}
  \item Prevents duplication and redundancy: Ensures that changes are only applied if necessary, making configuration management predictable and reliable.
\end{itemize}
\end{column}
\end{columns}
\end{frame}

% ==============================================================
% Slide 31 - Now we are ready
% ==============================================================
\begin{frame}[fragile]{Now we are ready}
\begin{itemize}
  \item Oh, right. First time SSHing, getting the fingerprint message.
  \item Also, we are trying to connect as \texttt{ansible\_user}
\end{itemize}
\vspace{0.2cm}
\begin{termsmall}
ansible_user@ansible:~/secure_devops_demo$ ansible-playbook ping-servers.yml

PLAY [Ping servers] *******************************************************
TASK [Gathering Facts] ****************************************************
fatal: [192.168.122.101]: UNREACHABLE! => {"changed": false, "msg": "Failed
  to connect to the host via ssh: ansible_user@192.168.122.101: Permission
  denied (publickey).", "unreachable": true}
The authenticity of host '192.168.122.102 (192.168.122.102)' can't be
  established.
ED25519 key fingerprint is SHA256:PzXwrkP3W+...
Are you sure you want to continue connecting (yes/no/[fingerprint])? yes
fatal: [192.168.122.102]: UNREACHABLE! => {"changed": false, "msg": "Failed
  to connect to the host via ssh: Permission denied (publickey)."}
\end{termsmall}
\end{frame}

% ==============================================================
% Slide 32 - Add ansible_user to servers
% ==============================================================
\begin{frame}[fragile]{Add \texttt{ansible\_user} to servers}
\begin{columns}[T]
\begin{column}{0.48\textwidth}
\begin{termsmall}
---
- name: Add ansible_user to the sudo
    group, set up SSH key, and set
    password
  hosts: servers
  become: true
  remote_user: root  vars_files:
    - vault.yml
  tasks:
    - name: Install passlib on remote
        hosts (if required)
      ansible.builtin.package:
        name: python3-passlib
        state: present
    - name: Ensure ansible_user is in
        the sudo group with a hashed
        password
      user:
        name: ansible_user
        groups: sudo
        append: true
        password: "{{ ansible_user_password
          | password_hash('sha512') }}"
\end{termsmall}
\end{column}
\begin{column}{0.48\textwidth}
\begin{termsmall}
ansible_user@ansible:~/secure_devops_demo$
  ansible-playbook add_ansible_user.yml

PLAY [Add ansible_user to the sudo
  group, set up SSH key, and set
  password] ************************

TASK [Gathering Facts] *************
ok: [192.168.122.101]
ok: [192.168.122.103]
ok: [192.168.122.102]

TASK [Install passlib...] **********
ok: [192.168.122.101]
ok: [192.168.122.102]
ok: [192.168.122.103]

TASK [Ensure ansible_user...] ******
changed: [192.168.122.101]
changed: [192.168.122.102]
changed: [192.168.122.103]

TASK [Create .ssh directory...] ****
ok: [192.168.122.101]
ok: [192.168.122.102]
ok: [192.168.122.103]

TASK [Add SSH public key...] *******
ok: [192.168.122.101]
ok: [192.168.122.102]
ok: [192.168.122.103]
\end{termsmall}
\end{column}
\end{columns}
\end{frame}

% ==============================================================
% Slide 33 - Wait, what?
% ==============================================================
\begin{frame}[fragile]{Switch SSH config to \texttt{ansible\_user}}
\begin{termsmall}
ansible_user@ansible:~/secure_devops_demo$ cat ~/.ssh/config
# BEGIN cyberforge-demo
Host server1
    HostName 192.168.122.101
    User ansible_user
    IdentityFile ~/.ssh/DemoSSHKey
    AddKeysToAgent yes
    StrictHostKeyChecking accept-new

Host server2
    HostName 192.168.122.102
    User ansible_user
    IdentityFile ~/.ssh/DemoSSHKey
    AddKeysToAgent yes
    StrictHostKeyChecking accept-new

Host server3
    HostName 192.168.122.103
    User ansible_user
    IdentityFile ~/.ssh/DemoSSHKey
    AddKeysToAgent yes
    StrictHostKeyChecking accept-new
# END cyberforge-demo
\end{termsmall}
\end{frame}

% ==============================================================
% Slide 34 - Verify ansible_user can connect
% ==============================================================
\begin{frame}[fragile]{Verify \texttt{ansible\_user} can connect}
\begin{term}
ansible_user@ansible:~/secure_devops_demo$ ansible-playbook ping-servers.yml

PLAY [Ping servers] ***************************************************

TASK [Gathering Facts] ************************************************
ok: [192.168.122.101]
ok: [192.168.122.102]
ok: [192.168.122.103]

TASK [Ping the localhost] *********************************************
ok: [192.168.122.101]
ok: [192.168.122.102]
ok: [192.168.122.103]

PLAY RECAP ************************************************************
192.168.122.101    : ok=2    changed=0    unreachable=0    failed=0
192.168.122.102    : ok=2    changed=0    unreachable=0    failed=0
192.168.122.103    : ok=2    changed=0    unreachable=0    failed=0
\end{term}
\vspace{0.2cm}
\textbf{Verify before you lock the door.}
\end{frame}

% ==============================================================
% Slide 35 - SSHD Hardening
% ==============================================================
\begin{frame}[fragile]{SSHD Hardening}
\begin{columns}[T]
\begin{column}{0.48\textwidth}
\begin{termsmall}
---
- name: Configure SSHD for Ubuntu
    22.04, 24.04, and 24.10
  hosts: servers
  become: true
  vars_files:    - vault.yml  tasks:
    - name: Ensure root login is
        disabled
      lineinfile:
        path: /etc/ssh/sshd_config
        regexp: '^#?PermitRootLogin'
        line: 'PermitRootLogin no'
        state: present
    - name: Ensure public key
        authentication is enabled
      lineinfile:
        path: /etc/ssh/sshd_config
        regexp: '^#?PubkeyAuthentication'
        line: 'PubkeyAuthentication yes'
        state: present
    - name: Disable password
        authentication
      lineinfile:
        path: /etc/ssh/sshd_config
        regexp: '^#?PasswordAuthentication'
        line: 'PasswordAuthentication no'
        state: present
    - name: Restart SSHD
      service:
        name: ssh
        state: restarted
\end{termsmall}
\end{column}
\begin{column}{0.48\textwidth}
\begin{termsmall}
ansible_user@ansible:~/secure_devops_demo$
  ansible-playbook
    sshd_hardening-servers.yml

PLAY [Configure SSHD for Ubuntu
  22.04, 24.04, and 24.10] *********

TASK [Gathering Facts] *************
ok: [192.168.122.101]
ok: [192.168.122.103]
ok: [192.168.122.102]

TASK [Ensure root login is
  disabled] **************************
changed: [192.168.122.101]
changed: [192.168.122.102]
changed: [192.168.122.103]

TASK [Ensure public key auth...] ***
changed: [192.168.122.101]
changed: [192.168.122.102]
changed: [192.168.122.103]

TASK [Disable password auth...] ****
changed: [192.168.122.101]
changed: [192.168.122.102]
changed: [192.168.122.103]

TASK [Restart SSHD...] *************
changed: [192.168.122.101]
changed: [192.168.122.102]
changed: [192.168.122.103]
\end{termsmall}
\end{column}
\end{columns}
\end{frame}

% ==============================================================
% Slide 36 - Now we can ping servers
% ==============================================================
\begin{frame}[fragile]{Now we can ping servers}
\begin{term}
ansible_user@ansible:~/secure_devops_demo$ ansible-playbook ping-servers.yml

PLAY [Ping servers] ***************************************************

TASK [Gathering Facts] ************************************************
ok: [192.168.122.101]
ok: [192.168.122.102]
ok: [192.168.122.103]

TASK [Ping the localhost] *********************************************
ok: [192.168.122.101]
ok: [192.168.122.102]
ok: [192.168.122.103]

PLAY RECAP ************************************************************
192.168.122.101    : ok=2    changed=0    unreachable=0    failed=0
192.168.122.102    : ok=2    changed=0    unreachable=0    failed=0
192.168.122.103    : ok=2    changed=0    unreachable=0    failed=0
\end{term}
\end{frame}

% ==============================================================
% Slide 37 - Ad-hoc commands
% ==============================================================
\begin{frame}[fragile]{Ad-hoc commands}
\begin{columns}[T]
\begin{column}{0.48\textwidth}
\texttt{ansible servers -m shell -a "df -h"}
\vspace{0.1cm}
\begin{termsmall}
192.168.122.101 | CHANGED | rc=0 >>
Filesystem  Size Used Avail Use% Mounted
/dev/vda1   9.6G 1.7G  7.9G  18% /
tmpfs        46M 1008K  45M   3% /run
...
192.168.122.102 | CHANGED | rc=0 >>
Filesystem  Size Used Avail Use% Mounted
/dev/vda1   9.6G 1.7G  7.9G  18% /
...
192.168.122.103 | CHANGED | rc=0 >>
Filesystem  Size Used Avail Use% Mounted
/dev/vda1   9.6G 1.7G  7.9G  18% /
\end{termsmall}
\end{column}
\begin{column}{0.48\textwidth}
\texttt{ansible servers -a "ss -tuln"}
\vspace{0.1cm}
\begin{termsmall}
192.168.122.101 | CHANGED | rc=0 >>
Netid State Recv-Q Send-Q
  Local Address:Port
udp   UNCONN 0      0
  127.0.0.53%lo:53
tcp   LISTEN 0   4096
  127.0.0.53%lo:53
tcp   LISTEN 0    128
  0.0.0.0:22
tcp   LISTEN 0    128
  [::]:22
192.168.122.102 | CHANGED | rc=0 >>
...
192.168.122.103 | CHANGED | rc=0 >>
...
\end{termsmall}
\end{column}
\end{columns}
\end{frame}

% ==============================================================
% Slide 38 - Summary
% ==============================================================
\begin{frame}{Summary}
\begin{itemize}
  \item DevOps is awesome!
  \item Many demos and tutorials skip over how to set up a secure environment for Ansible
  \begin{itemize}
    \item SSH \& GPG Keys
    \item SSH-Agent
    \item GPG-Agent
    \item \texttt{pass} -- The Unix Password Manager
    \item Ansible Vault
  \end{itemize}
\end{itemize}
\end{frame}

% ==============================================================
% Slide 39 - Questions?
% ==============================================================
\begin{frame}[plain]
\vspace{3cm}
\begin{beamercolorbox}[wd=\paperwidth,ht=2.5cm,dp=0.8cm]{frametitle}
  \hfill{\Huge Questions?}\hspace{1cm}
  \vspace{0.5cm}

  \hfill{\small CONTACT:}\hspace{1cm}

  \hfill{\large\textbf{jkirch (at) tcc dot edu}}\hspace{1cm}
\end{beamercolorbox}
\vfill
\begin{center}
{\tiny\textcopyright\ CC BY-NC-ND --- This license enables reusers to copy and distribute the material in any medium or format in unadapted form only, for noncommercial purposes only, and only so long as attribution is given to the creator.}
\end{center}
\end{frame}

\end{document}
